% !TEX program = xelatex
% ============================================================
%   控制理论笔记 —— 主文档
%
%   编译：在 nvim 里对 main.tex 按 \ll（vimtex 连续编译，latexmk + XeLaTeX）
%   输出布局（对应 ~/.config/nvim/lua/plugins/latex.lua）：
%       辅助文件(.aux/.log/.toc/.synctex) → build/
%       PDF(main.pdf)                     → 本目录（根目录）
%       \lv 正向跳转看 PDF，zathura 反向跳转回源码
%
%   新增章节：chapters/ 下建 NN_名字.tex（首行 % !TEX root = ../main.tex），
%             然后在下方正文区 \include 一行。
% ============================================================

\documentclass[a4paper, 11pt, fontset=none, oneside, openany]{ctexbook}

% ---------- 导言区：整体风格（字体/配色/页面/框，见 preamble.tex） ----------
% ============================================================
%   导言区 / 序言（preamble）—— 管整体风格
%   由 main.tex 用  % ============================================================
%   导言区 / 序言（preamble）—— 管整体风格
%   由 main.tex 用  % ============================================================
%   导言区 / 序言（preamble）—— 管整体风格
%   由 main.tex 用  \input{preamble}  引入（在 \begin{document} 之前）。
%   统一改风格（字体、配色、页面、框）都是改这里。
% ============================================================

% ---------- 字体（XeLaTeX，已由 ctex 加载 xeCJK/fontspec） ----------
% 正文思源宋体；中文「斜体」按国内惯例映射成楷体（\textit / \emph 即楷体）
\setCJKmainfont[ItalicFont=LXGW WenKai, BoldItalicFont=LXGW WenKai]{Noto Serif CJK SC}
\setCJKsansfont{Noto Sans CJK SC}    % 标题中文：思源黑体
\setCJKmonofont{LXGW WenKai Mono}    % 中文等宽：霞鹜文楷 Mono
\setmonofont{JetBrainsMono NF}       % 拉丁等宽（代码）：JetBrains Mono Nerd Font
% 拉丁正文与数学保留 Computer Modern。

% ---------- 数学 ----------
\usepackage{amsmath,amssymb}
\usepackage{bm}
\usepackage{mathtools}
\numberwithin{equation}{chapter}   % 公式编号按章：1.1 / 1.2 ...
\numberwithin{figure}{chapter}     % 图编号按章

% ---------- 页面（A4） ----------
\usepackage[a4paper, top=2.3cm, bottom=2.3cm, left=2.2cm, right=2.2cm,
  headheight=14pt, headsep=6mm, footskip=11mm]{geometry}
\usepackage{setspace}
\onehalfspacing

% ---------- 图形 ----------
\usepackage{graphicx}
\usepackage{tikz}
\usetikzlibrary{calc, arrows.meta, positioning, shapes.geometric}

% ---------- 表格 ----------
\usepackage{booktabs}
\usepackage{array}

% ---------- 物理单位 ----------
\usepackage{siunitx}

% ---------- 列表间距 ----------
\usepackage{enumitem}
\setlist{topsep=3pt, partopsep=0pt, itemsep=2pt, parsep=0pt}

% ---------- 颜色（知识块调色板 + 封面） ----------
\usepackage{xcolor}
\definecolor{bl}{HTML}{3B6BB0}   % 蓝 — 定义
\definecolor{tl}{HTML}{1F9B8E}   % 青 — 理解
\definecolor{gr}{HTML}{4C9F70}   % 绿 — 例
\definecolor{og}{HTML}{E8933A}   % 橙 — 证明思路
\definecolor{pu}{HTML}{7B5EA7}   % 紫 — 命题
\definecolor{rd}{HTML}{C0504D}   % 红 — 定理
\definecolor{in}{HTML}{283593}   % 靛 — 引理
\definecolor{pk}{HTML}{C86B98}   % 粉 — 推论
\definecolor{gy}{HTML}{6B7280}   % 灰 — 注
\definecolor{coverblue}{HTML}{283593}   % 封面深蓝
\definecolor{coverlight}{HTML}{5C6BC0}  % 封面浅蓝
\definecolor{covergray}{HTML}{5A5F6E}   % 封面灰
\definecolor{link}{HTML}{1F6F8B}

% ---------- 页眉页脚 ----------
\usepackage{fancyhdr}
\pagestyle{fancy}
\fancyhf{}
\fancyhead[L]{\small\nouppercase{\rightmark}}
\fancyhead[R]{\small\thepage}
\renewcommand{\headrulewidth}{0.4pt}

% ---------- 章节标题（加大醒目） ----------
\usepackage{titlesec}
\titleformat{\chapter}{\centering\fontsize{20}{26}\selectfont\bfseries}{\Huge 第\chinese{chapter}章}{0.6em}{}
\titlespacing*{\chapter}{0pt}{-16pt}{30pt}
\titleformat{\section}{\fontsize{16}{20}\selectfont\bfseries}{\thesection}{0.6em}{}
\titlespacing*{\section}{0pt}{16pt}{6pt}
\titleformat{\subsection}{\fontsize{14}{18}\selectfont\bfseries}{\thesubsection}{0.6em}{}
\titlespacing*{\subsection}{0pt}{12pt}{4pt}
\titleformat{\subsubsection}{\fontsize{12}{15}\selectfont\bfseries}{\thesubsubsection}{0.6em}{}

% ---------- 目录（加大字号） ----------
\usepackage{tocloft}
\renewcommand{\cfttoctitlefont}{\Huge\bfseries}
\renewcommand{\cftbeforetoctitleskip}{2em}
\renewcommand{\cftaftertoctitleskip}{1.2em}
\renewcommand{\cftchapfont}{\large\bfseries}
\renewcommand{\cftchappagefont}{\large}
\renewcommand{\cftsecfont}{\normalsize}
\renewcommand{\cftsubsecfont}{\normalsize}
\renewcommand{\cftbeforechapskip}{1.2em}
\renewcommand{\cftbeforesecskip}{0.4em}

% ---------- 超链接 ----------
\usepackage[colorlinks=true,linkcolor=link,citecolor=link,urlcolor=link]{hyperref}

% ---------- 图形/表格题注 ----------
\usepackage{caption}
\captionsetup{font=small, labelfont=bf}

% ---------- 知识块（饱和色标题条 + 同色左边线 + 极淡底 + 方角） ----------
\usepackage[most]{tcolorbox}
\tcbset{
  jbox/.style = {
    enhanced, breakable, arc=3pt,
    boxrule=0.6pt, leftrule=3pt,   % 粗左边线（同标题色）
    left=10pt, right=10pt, top=8pt, bottom=8pt,
    before skip=12pt plus 2pt, after skip=12pt plus 2pt,
    fonttitle=\bfseries, coltitle=white,
    toptitle=1.5pt, bottomtitle=1.5pt,
    theorem style=standard,
    separator sign none, terminator sign none,
    description delimiters={\kern-0.42em（}{）},
    % parbox 用默认(true)：盒内文字自动按盒宽换行，不会溢出边框
  },
}
% 彩色定理环境：一行一个，颜色即第 3 个参数
%   用法：\newcoloredtheorem{环境名}{显示名}{颜色}{标签前缀}
\newcommand{\newcoloredtheorem}[4]{%
  \newtcbtheorem[number within=chapter]{#1}{#2}{jbox, colback=#3!8, colframe=#3, colbacktitle=#3}{#4}%
}
\newcoloredtheorem{definition}{定义}{bl}{def}
\newcoloredtheorem{theorem}{定理}{rd}{thm}
\newcoloredtheorem{lemma}{引理}{in}{lem}
\newcoloredtheorem{proposition}{命题}{pu}{prop}
\newcoloredtheorem{corollary}{推论}{pk}{cor}
\newcoloredtheorem{example}{例}{gr}{ex}
\newcoloredtheorem{remark}{注}{gy}{rem}
% 通用彩色信息块：#1=颜色(可选，默认蓝) #2=标题(必填)
%   用法：\begin{notebox}[颜色]{标题} ... \end{notebox}
\newtcolorbox{notebox}[2][bl]{jbox, colback=#1!8, colframe=#1, colbacktitle=#1, title={#2}}

% ---------- 代码框（语法高亮 + 代码框） ----------
% 语法高亮用 listings，外面套 tcolorbox 做个带标题的代码框。
%   用法：\begin{codebox}[语言]{标题}  代码  \end{codebox}
\usepackage{listings}
\definecolor{codebg}{HTML}{F5F6F8}   % 代码区背景
\definecolor{codeframe}{HTML}{9AA2B0} % 代码框边框
\definecolor{kw}{HTML}{4A6FD0}        % 关键字
\definecolor{st}{HTML}{4C9F70}        % 字符串
\definecolor{cm}{HTML}{9AA0A6}        % 注释
\definecolor{nm}{HTML}{E8933A}        % 行号
\lstset{
  basicstyle=\ttfamily\small,
  keywordstyle=\color{kw}\bfseries,
  stringstyle=\color{st},
  commentstyle=\color{cm}\itshape,
  numberstyle=\color{nm}\tiny,
  numbers=left, numbersep=6pt,
  showstringspaces=false, breaklines=true, tabsize=2,
  frame=none,
}
% listings 不带内置的 latex 语言，自定义一个；常见命令可高亮
\lstdefinelanguage{latex}{%
  morekeywords={documentclass,usepackage,begin,end,chapter,section,subsection,subsubsection,
    label,ref,eqref,caption,include,input,equation,align,figure,tabular,hline,centering,
    includegraphics,textbf,textit,emph,texttt,maketitle,tableofcontents,titlepage,body},
  morecomment=[l]{\%},
  sensitive=true,
}
\lstalias{tex}{latex}
\lstalias{sh}{bash}
% 无高亮的纯文本（目录/说明用），避免把 ​# 当注释导致解析问题
\lstdefinelanguage{text}{sensitive=true}
\newtcblisting{codebox}[2][bash]{%
  enhanced, breakable, listing only,
  listing options={language=#1},
  colback=codebg, colframe=codeframe, leftrule=3pt,
  colbacktitle=bl!80!black, coltitle=white, fonttitle=\bfseries,
  arc=2mm, boxrule=0.5pt,
  left=6pt, right=6pt, top=4pt, bottom=4pt,
  title={#2},
}
  引入（在 \begin{document} 之前）。
%   统一改风格（字体、配色、页面、框）都是改这里。
% ============================================================

% ---------- 字体（XeLaTeX，已由 ctex 加载 xeCJK/fontspec） ----------
% 正文思源宋体；中文「斜体」按国内惯例映射成楷体（\textit / \emph 即楷体）
\setCJKmainfont[ItalicFont=LXGW WenKai, BoldItalicFont=LXGW WenKai]{Noto Serif CJK SC}
\setCJKsansfont{Noto Sans CJK SC}    % 标题中文：思源黑体
\setCJKmonofont{LXGW WenKai Mono}    % 中文等宽：霞鹜文楷 Mono
\setmonofont{JetBrainsMono NF}       % 拉丁等宽（代码）：JetBrains Mono Nerd Font
% 拉丁正文与数学保留 Computer Modern。

% ---------- 数学 ----------
\usepackage{amsmath,amssymb}
\usepackage{bm}
\usepackage{mathtools}
\numberwithin{equation}{chapter}   % 公式编号按章：1.1 / 1.2 ...
\numberwithin{figure}{chapter}     % 图编号按章

% ---------- 页面（A4） ----------
\usepackage[a4paper, top=2.3cm, bottom=2.3cm, left=2.2cm, right=2.2cm,
  headheight=14pt, headsep=6mm, footskip=11mm]{geometry}
\usepackage{setspace}
\onehalfspacing

% ---------- 图形 ----------
\usepackage{graphicx}
\usepackage{tikz}
\usetikzlibrary{calc, arrows.meta, positioning, shapes.geometric}

% ---------- 表格 ----------
\usepackage{booktabs}
\usepackage{array}

% ---------- 物理单位 ----------
\usepackage{siunitx}

% ---------- 列表间距 ----------
\usepackage{enumitem}
\setlist{topsep=3pt, partopsep=0pt, itemsep=2pt, parsep=0pt}

% ---------- 颜色（知识块调色板 + 封面） ----------
\usepackage{xcolor}
\definecolor{bl}{HTML}{3B6BB0}   % 蓝 — 定义
\definecolor{tl}{HTML}{1F9B8E}   % 青 — 理解
\definecolor{gr}{HTML}{4C9F70}   % 绿 — 例
\definecolor{og}{HTML}{E8933A}   % 橙 — 证明思路
\definecolor{pu}{HTML}{7B5EA7}   % 紫 — 命题
\definecolor{rd}{HTML}{C0504D}   % 红 — 定理
\definecolor{in}{HTML}{283593}   % 靛 — 引理
\definecolor{pk}{HTML}{C86B98}   % 粉 — 推论
\definecolor{gy}{HTML}{6B7280}   % 灰 — 注
\definecolor{coverblue}{HTML}{283593}   % 封面深蓝
\definecolor{coverlight}{HTML}{5C6BC0}  % 封面浅蓝
\definecolor{covergray}{HTML}{5A5F6E}   % 封面灰
\definecolor{link}{HTML}{1F6F8B}

% ---------- 页眉页脚 ----------
\usepackage{fancyhdr}
\pagestyle{fancy}
\fancyhf{}
\fancyhead[L]{\small\nouppercase{\rightmark}}
\fancyhead[R]{\small\thepage}
\renewcommand{\headrulewidth}{0.4pt}

% ---------- 章节标题（加大醒目） ----------
\usepackage{titlesec}
\titleformat{\chapter}{\centering\fontsize{20}{26}\selectfont\bfseries}{\Huge 第\chinese{chapter}章}{0.6em}{}
\titlespacing*{\chapter}{0pt}{-16pt}{30pt}
\titleformat{\section}{\fontsize{16}{20}\selectfont\bfseries}{\thesection}{0.6em}{}
\titlespacing*{\section}{0pt}{16pt}{6pt}
\titleformat{\subsection}{\fontsize{14}{18}\selectfont\bfseries}{\thesubsection}{0.6em}{}
\titlespacing*{\subsection}{0pt}{12pt}{4pt}
\titleformat{\subsubsection}{\fontsize{12}{15}\selectfont\bfseries}{\thesubsubsection}{0.6em}{}

% ---------- 目录（加大字号） ----------
\usepackage{tocloft}
\renewcommand{\cfttoctitlefont}{\Huge\bfseries}
\renewcommand{\cftbeforetoctitleskip}{2em}
\renewcommand{\cftaftertoctitleskip}{1.2em}
\renewcommand{\cftchapfont}{\large\bfseries}
\renewcommand{\cftchappagefont}{\large}
\renewcommand{\cftsecfont}{\normalsize}
\renewcommand{\cftsubsecfont}{\normalsize}
\renewcommand{\cftbeforechapskip}{1.2em}
\renewcommand{\cftbeforesecskip}{0.4em}

% ---------- 超链接 ----------
\usepackage[colorlinks=true,linkcolor=link,citecolor=link,urlcolor=link]{hyperref}

% ---------- 图形/表格题注 ----------
\usepackage{caption}
\captionsetup{font=small, labelfont=bf}

% ---------- 知识块（饱和色标题条 + 同色左边线 + 极淡底 + 方角） ----------
\usepackage[most]{tcolorbox}
\tcbset{
  jbox/.style = {
    enhanced, breakable, arc=3pt,
    boxrule=0.6pt, leftrule=3pt,   % 粗左边线（同标题色）
    left=10pt, right=10pt, top=8pt, bottom=8pt,
    before skip=12pt plus 2pt, after skip=12pt plus 2pt,
    fonttitle=\bfseries, coltitle=white,
    toptitle=1.5pt, bottomtitle=1.5pt,
    theorem style=standard,
    separator sign none, terminator sign none,
    description delimiters={\kern-0.42em（}{）},
    % parbox 用默认(true)：盒内文字自动按盒宽换行，不会溢出边框
  },
}
% 彩色定理环境：一行一个，颜色即第 3 个参数
%   用法：\newcoloredtheorem{环境名}{显示名}{颜色}{标签前缀}
\newcommand{\newcoloredtheorem}[4]{%
  \newtcbtheorem[number within=chapter]{#1}{#2}{jbox, colback=#3!8, colframe=#3, colbacktitle=#3}{#4}%
}
\newcoloredtheorem{definition}{定义}{bl}{def}
\newcoloredtheorem{theorem}{定理}{rd}{thm}
\newcoloredtheorem{lemma}{引理}{in}{lem}
\newcoloredtheorem{proposition}{命题}{pu}{prop}
\newcoloredtheorem{corollary}{推论}{pk}{cor}
\newcoloredtheorem{example}{例}{gr}{ex}
\newcoloredtheorem{remark}{注}{gy}{rem}
% 通用彩色信息块：#1=颜色(可选，默认蓝) #2=标题(必填)
%   用法：\begin{notebox}[颜色]{标题} ... \end{notebox}
\newtcolorbox{notebox}[2][bl]{jbox, colback=#1!8, colframe=#1, colbacktitle=#1, title={#2}}

% ---------- 代码框（语法高亮 + 代码框） ----------
% 语法高亮用 listings，外面套 tcolorbox 做个带标题的代码框。
%   用法：\begin{codebox}[语言]{标题}  代码  \end{codebox}
\usepackage{listings}
\definecolor{codebg}{HTML}{F5F6F8}   % 代码区背景
\definecolor{codeframe}{HTML}{9AA2B0} % 代码框边框
\definecolor{kw}{HTML}{4A6FD0}        % 关键字
\definecolor{st}{HTML}{4C9F70}        % 字符串
\definecolor{cm}{HTML}{9AA0A6}        % 注释
\definecolor{nm}{HTML}{E8933A}        % 行号
\lstset{
  basicstyle=\ttfamily\small,
  keywordstyle=\color{kw}\bfseries,
  stringstyle=\color{st},
  commentstyle=\color{cm}\itshape,
  numberstyle=\color{nm}\tiny,
  numbers=left, numbersep=6pt,
  showstringspaces=false, breaklines=true, tabsize=2,
  frame=none,
}
% listings 不带内置的 latex 语言，自定义一个；常见命令可高亮
\lstdefinelanguage{latex}{%
  morekeywords={documentclass,usepackage,begin,end,chapter,section,subsection,subsubsection,
    label,ref,eqref,caption,include,input,equation,align,figure,tabular,hline,centering,
    includegraphics,textbf,textit,emph,texttt,maketitle,tableofcontents,titlepage,body},
  morecomment=[l]{\%},
  sensitive=true,
}
\lstalias{tex}{latex}
\lstalias{sh}{bash}
% 无高亮的纯文本（目录/说明用），避免把 ​# 当注释导致解析问题
\lstdefinelanguage{text}{sensitive=true}
\newtcblisting{codebox}[2][bash]{%
  enhanced, breakable, listing only,
  listing options={language=#1},
  colback=codebg, colframe=codeframe, leftrule=3pt,
  colbacktitle=bl!80!black, coltitle=white, fonttitle=\bfseries,
  arc=2mm, boxrule=0.5pt,
  left=6pt, right=6pt, top=4pt, bottom=4pt,
  title={#2},
}
  引入（在 \begin{document} 之前）。
%   统一改风格（字体、配色、页面、框）都是改这里。
% ============================================================

% ---------- 字体（XeLaTeX，已由 ctex 加载 xeCJK/fontspec） ----------
% 正文思源宋体；中文「斜体」按国内惯例映射成楷体（\textit / \emph 即楷体）
\setCJKmainfont[ItalicFont=LXGW WenKai, BoldItalicFont=LXGW WenKai]{Noto Serif CJK SC}
\setCJKsansfont{Noto Sans CJK SC}    % 标题中文：思源黑体
\setCJKmonofont{LXGW WenKai Mono}    % 中文等宽：霞鹜文楷 Mono
\setmonofont{JetBrainsMono NF}       % 拉丁等宽（代码）：JetBrains Mono Nerd Font
% 拉丁正文与数学保留 Computer Modern。

% ---------- 数学 ----------
\usepackage{amsmath,amssymb}
\usepackage{bm}
\usepackage{mathtools}
\numberwithin{equation}{chapter}   % 公式编号按章：1.1 / 1.2 ...
\numberwithin{figure}{chapter}     % 图编号按章

% ---------- 页面（A4） ----------
\usepackage[a4paper, top=2.3cm, bottom=2.3cm, left=2.2cm, right=2.2cm,
  headheight=14pt, headsep=6mm, footskip=11mm]{geometry}
\usepackage{setspace}
\onehalfspacing

% ---------- 图形 ----------
\usepackage{graphicx}
\usepackage{tikz}
\usetikzlibrary{calc, arrows.meta, positioning, shapes.geometric}

% ---------- 表格 ----------
\usepackage{booktabs}
\usepackage{array}

% ---------- 物理单位 ----------
\usepackage{siunitx}

% ---------- 列表间距 ----------
\usepackage{enumitem}
\setlist{topsep=3pt, partopsep=0pt, itemsep=2pt, parsep=0pt}

% ---------- 颜色（知识块调色板 + 封面） ----------
\usepackage{xcolor}
\definecolor{bl}{HTML}{3B6BB0}   % 蓝 — 定义
\definecolor{tl}{HTML}{1F9B8E}   % 青 — 理解
\definecolor{gr}{HTML}{4C9F70}   % 绿 — 例
\definecolor{og}{HTML}{E8933A}   % 橙 — 证明思路
\definecolor{pu}{HTML}{7B5EA7}   % 紫 — 命题
\definecolor{rd}{HTML}{C0504D}   % 红 — 定理
\definecolor{in}{HTML}{283593}   % 靛 — 引理
\definecolor{pk}{HTML}{C86B98}   % 粉 — 推论
\definecolor{gy}{HTML}{6B7280}   % 灰 — 注
\definecolor{coverblue}{HTML}{283593}   % 封面深蓝
\definecolor{coverlight}{HTML}{5C6BC0}  % 封面浅蓝
\definecolor{covergray}{HTML}{5A5F6E}   % 封面灰
\definecolor{link}{HTML}{1F6F8B}

% ---------- 页眉页脚 ----------
\usepackage{fancyhdr}
\pagestyle{fancy}
\fancyhf{}
\fancyhead[L]{\small\nouppercase{\rightmark}}
\fancyhead[R]{\small\thepage}
\renewcommand{\headrulewidth}{0.4pt}

% ---------- 章节标题（加大醒目） ----------
\usepackage{titlesec}
\titleformat{\chapter}{\centering\fontsize{20}{26}\selectfont\bfseries}{\Huge 第\chinese{chapter}章}{0.6em}{}
\titlespacing*{\chapter}{0pt}{-16pt}{30pt}
\titleformat{\section}{\fontsize{16}{20}\selectfont\bfseries}{\thesection}{0.6em}{}
\titlespacing*{\section}{0pt}{16pt}{6pt}
\titleformat{\subsection}{\fontsize{14}{18}\selectfont\bfseries}{\thesubsection}{0.6em}{}
\titlespacing*{\subsection}{0pt}{12pt}{4pt}
\titleformat{\subsubsection}{\fontsize{12}{15}\selectfont\bfseries}{\thesubsubsection}{0.6em}{}

% ---------- 目录（加大字号） ----------
\usepackage{tocloft}
\renewcommand{\cfttoctitlefont}{\Huge\bfseries}
\renewcommand{\cftbeforetoctitleskip}{2em}
\renewcommand{\cftaftertoctitleskip}{1.2em}
\renewcommand{\cftchapfont}{\large\bfseries}
\renewcommand{\cftchappagefont}{\large}
\renewcommand{\cftsecfont}{\normalsize}
\renewcommand{\cftsubsecfont}{\normalsize}
\renewcommand{\cftbeforechapskip}{1.2em}
\renewcommand{\cftbeforesecskip}{0.4em}

% ---------- 超链接 ----------
\usepackage[colorlinks=true,linkcolor=link,citecolor=link,urlcolor=link]{hyperref}

% ---------- 图形/表格题注 ----------
\usepackage{caption}
\captionsetup{font=small, labelfont=bf}

% ---------- 知识块（饱和色标题条 + 同色左边线 + 极淡底 + 方角） ----------
\usepackage[most]{tcolorbox}
\tcbset{
  jbox/.style = {
    enhanced, breakable, arc=3pt,
    boxrule=0.6pt, leftrule=3pt,   % 粗左边线（同标题色）
    left=10pt, right=10pt, top=8pt, bottom=8pt,
    before skip=12pt plus 2pt, after skip=12pt plus 2pt,
    fonttitle=\bfseries, coltitle=white,
    toptitle=1.5pt, bottomtitle=1.5pt,
    theorem style=standard,
    separator sign none, terminator sign none,
    description delimiters={\kern-0.42em（}{）},
    % parbox 用默认(true)：盒内文字自动按盒宽换行，不会溢出边框
  },
}
% 彩色定理环境：一行一个，颜色即第 3 个参数
%   用法：\newcoloredtheorem{环境名}{显示名}{颜色}{标签前缀}
\newcommand{\newcoloredtheorem}[4]{%
  \newtcbtheorem[number within=chapter]{#1}{#2}{jbox, colback=#3!8, colframe=#3, colbacktitle=#3}{#4}%
}
\newcoloredtheorem{definition}{定义}{bl}{def}
\newcoloredtheorem{theorem}{定理}{rd}{thm}
\newcoloredtheorem{lemma}{引理}{in}{lem}
\newcoloredtheorem{proposition}{命题}{pu}{prop}
\newcoloredtheorem{corollary}{推论}{pk}{cor}
\newcoloredtheorem{example}{例}{gr}{ex}
\newcoloredtheorem{remark}{注}{gy}{rem}
% 通用彩色信息块：#1=颜色(可选，默认蓝) #2=标题(必填)
%   用法：\begin{notebox}[颜色]{标题} ... \end{notebox}
\newtcolorbox{notebox}[2][bl]{jbox, colback=#1!8, colframe=#1, colbacktitle=#1, title={#2}}

% ---------- 代码框（语法高亮 + 代码框） ----------
% 语法高亮用 listings，外面套 tcolorbox 做个带标题的代码框。
%   用法：\begin{codebox}[语言]{标题}  代码  \end{codebox}
\usepackage{listings}
\definecolor{codebg}{HTML}{F5F6F8}   % 代码区背景
\definecolor{codeframe}{HTML}{9AA2B0} % 代码框边框
\definecolor{kw}{HTML}{4A6FD0}        % 关键字
\definecolor{st}{HTML}{4C9F70}        % 字符串
\definecolor{cm}{HTML}{9AA0A6}        % 注释
\definecolor{nm}{HTML}{E8933A}        % 行号
\lstset{
  basicstyle=\ttfamily\small,
  keywordstyle=\color{kw}\bfseries,
  stringstyle=\color{st},
  commentstyle=\color{cm}\itshape,
  numberstyle=\color{nm}\tiny,
  numbers=left, numbersep=6pt,
  showstringspaces=false, breaklines=true, tabsize=2,
  frame=none,
}
% listings 不带内置的 latex 语言，自定义一个；常见命令可高亮
\lstdefinelanguage{latex}{%
  morekeywords={documentclass,usepackage,begin,end,chapter,section,subsection,subsubsection,
    label,ref,eqref,caption,include,input,equation,align,figure,tabular,hline,centering,
    includegraphics,textbf,textit,emph,texttt,maketitle,tableofcontents,titlepage,body},
  morecomment=[l]{\%},
  sensitive=true,
}
\lstalias{tex}{latex}
\lstalias{sh}{bash}
% 无高亮的纯文本（目录/说明用），避免把 ​# 当注释导致解析问题
\lstdefinelanguage{text}{sensitive=true}
\newtcblisting{codebox}[2][bash]{%
  enhanced, breakable, listing only,
  listing options={language=#1},
  colback=codebg, colframe=codeframe, leftrule=3pt,
  colbacktitle=bl!80!black, coltitle=white, fonttitle=\bfseries,
  arc=2mm, boxrule=0.5pt,
  left=6pt, right=6pt, top=4pt, bottom=4pt,
  title={#2},
}



\begin{document}

% ---------- 封面（左侧色条 + 右上色块 + 大标题） ----------
\begin{titlepage}
\begin{tikzpicture}[remember picture, overlay]
  % 左竖条
  \fill[coverblue] ([xshift=0pt]current page.north west) rectangle ([xshift=1.6cm]current page.south west);
  % 右上两个小色块
  \fill[coverlight] ([xshift=-2.6cm,yshift=-1.2cm]current page.north east) rectangle ([xshift=-1.6cm,yshift=-0.2cm]current page.north east);
  \fill[coverblue] ([xshift=-1.6cm,yshift=-1.2cm]current page.north east) rectangle ([xshift=-0.6cm,yshift=-0.2cm]current page.north east);
  % 文字：以页面左上角为原点，x=2.8cm，y 向下（yshift 为负）
  \node[anchor=west] at ([xshift=2.8cm,yshift=-6.2cm]current page.north west) {\color{covergray}\large 现代控制理论 \quad $\cdot$ \quad 个人学习笔记};
  \node[anchor=west] at ([xshift=2.8cm,yshift=-7.4cm]current page.north west) {\color{covergray}\rule{5.5cm}{0.4pt}};
  \node[anchor=west] at ([xshift=2.8cm,yshift=-10.8cm]current page.north west) {\color{coverblue}\fontsize{52}{60}\selectfont\bfseries\sffamily 控制理论};
  \node[anchor=west] at ([xshift=2.8cm,yshift=-14.4cm]current page.north west) {\color{coverblue}\fontsize{52}{60}\selectfont\bfseries\sffamily 与状态空间};
  \node[anchor=west] at ([xshift=2.8cm,yshift=-16.8cm]current page.north west) {\color{covergray}\rule{15.0cm}{0.6pt}};
  \node[anchor=west] at ([xshift=2.8cm,yshift=-22.0cm]current page.north west) {\color{coverblue}\Large\bfseries\sffamily Basstt \quad 著};
  \node[anchor=west] at ([xshift=2.8cm,yshift=-23.4cm]current page.north west) {\color{covergray}\normalsize\href{https://github.com/BassttElSevic}{github.com/BassttElSevic}};
  \node[anchor=west] at ([xshift=2.8cm,yshift=-24.8cm]current page.north west) {\color{covergray}\normalsize\href{https://github.com/BassttElSevic/MyNoteOfControlTheory}{github.com/BassttElSevic/MyNoteOfControlTheory}};
\end{tikzpicture}
\end{titlepage}

% ---------- 前言 ----------
\frontmatter
\tableofcontents

% ---------- 正文 ----------
\mainmatter
% 教程章显示为「第零章」：把章节计数从 0 开始，下一章才是「第一章」
\setcounter{chapter}{-1}
% !TEX root = ../main.tex

% ============================================================
%   第零章：怎么用这套模板记笔记
%   写给第一次接触 LaTeX 的人。每节都先「展示效果」再「给出源码」。
% ============================================================

\chapter{怎么用这套模板记笔记}
\label{ch:start}

这套笔记用 \LaTeX 写，编译成 PDF。这一章把「从头到尾怎么用」讲清楚。\textbf{每一节都分两部分}：先给效果（这一章里真实渲染出来的样子），再给能复制出这个效果的源码。

\section{先理解：源文件与成品}

你在 \texttt{.tex} 里写的不是最终样子，而是「怎么排版」的指令。写完要编译，才得到 \texttt{main.pdf}。

\begin{notebox}[bl]{核心循环}
改 \texttt{.tex} → 编译（\texttt{\textbackslash ll}）→ 看 \texttt{main.pdf}（\texttt{\textbackslash lv}）。
\end{notebox}

\section{编译与查看（\texttt{\textbackslash ll}）}

\subsection{写笔记的动作只有三步}

\begin{enumerate}
  \item 在 nvim 里打开 \texttt{main.tex}，按 \texttt{\textbackslash ll} 编译。\texttt{\textbackslash ll} 是连续编译：你保存一次它自动重编一次。
  \item 按 \texttt{\textbackslash lv} 打开 \texttt{main.pdf}（Zathura），光标跳到你源码所在处。
  \item 在 Zathura 里 \texttt{Ctrl + 左键} 点 PDF，跳回源码对应行（反向跳转）。
\end{enumerate}

\subsection{按 \texttt{\textbackslash ll} 之后发生了什么}

\begin{itemize}
  \item 辅助文件（\texttt{.aux}、\texttt{.log}、\texttt{.toc}、\texttt{.synctex}）写进 \texttt{build/}。
  \item \texttt{main.pdf} 写在项目根目录，紧挨 \texttt{main.tex}。
  \item 目录/引用要靠 \textbf{编两遍} 才正确。
\end{itemize}

\section{vimtex 快捷键与这套配置改了啥}

\subsection{vimtex 的按键}

\texttt{vimtex} 用 \texttt{<localleader>}（这里是 \textbackslash）开头的按键管编译。常用的几个：

\begin{table}[htbp]
  \centering
  \begin{tabular}{ll}
    \toprule
    按键 & 作用 \\
    \midrule
    \texttt{\textbackslash ll} & 开始/停止连续编译（latexmk）\\
    \texttt{\textbackslash lv} & 打开 PDF 并跳到光标处（zathura）\\
    \texttt{\textbackslash lc} & 清理辅助文件 \\
    \texttt{\textbackslash le} & 查看编译错误 \\
    \texttt{\textbackslash lt} & 章节/目录 \\
    \texttt{\textbackslash li} & 编译信息 \\
    \bottomrule
  \end{tabular}
\end{table}

最常用就是 \texttt{\textbackslash ll}（编译）和 \texttt{\textbackslash lv}（看 PDF）。

\subsection{这套配置在 nvim 里改了什么}

配置文件在 \texttt{\textasciitilde/.config/nvim/lua/plugins/latex.lua}，主要几处：

\begin{itemize}
  \item \textbf{编译器用 \texttt{latexmk}}：自动跑多遍、处理 \texttt{\textbackslash bibliography} 等。
  \item \textbf{默认引擎是 \texttt{xelatex}}：为了中文。
        \texttt{pdflatex} 处理不了中文字符；用 \texttt{xelatex} 才有 \texttt{fontspec} 字库。
        想单独某个文件用 pdflatex，在文件头写 \texttt{\% !TEX program = pdflatex}。
  \item \textbf{辅助文件进 \texttt{build/}}：\texttt{aux\_dir = "build"}，源码目录只留 \texttt{.tex} \textbf{和} \texttt{main.pdf}。
  \item \textbf{PDF 留在源码旁}：\texttt{out\_dir = ""}，方便 \texttt{\textbackslash lv}。
  \item \textbf{阅读器用 \texttt{zathura\_simple}}：支持 SyncTeX 正反向跳转，且不依赖 X11 的 \texttt{xdotool}（本机是 Wayland）。
  \item \textbf{补全和诊断交给 \texttt{texlab}}（LSP）：关闭了 vimtex 自己的补全（\texttt{vimtex\_complete\_enabled=0}）和诊断队列，避免重复。
  \item \textbf{texlab 编译手动触发}：\texttt{onSave = false}，不每次保存都编（大文档慢），要编用 \texttt{:LspTexlabBuild}。
  \item \textbf{chktex 静态检查}：打开和保存时跑，在你写的时候顺手查语法/标点问题。
\end{itemize}

\begin{notebox}[og]{想改这些}
这些都是 \texttt{~/.config/nvim/lua/plugins/latex.lua} 里的 \texttt{vim.g.vimtex\_* } 变量，改那一个文件即可，跟笔记 \texttt{.tex} 无关。
\end{notebox}

\section{这套模板的文件}

\begin{codebox}[text]{项目结构}
MyNoteOfControlTheory/
├── main.tex        # 主文档：文档类、封面、\include 各章节
├── preamble.tex    # 导言区：字体、配色、页面、知识块宏、代码框宏
├── chapters/       # 各章节 .tex
│   └── 00_getting_started.tex
├── images/         # 外部图片
├── build/          # 辅助文件（自动生成，已 gitignore）
├── main.pdf        # 编译成品（根目录）
├── .gitignore
└── README.md
\end{codebox}

\textbf{你平时只动 \texttt{main.tex} 和 \texttt{chapters/}。} 全局样式（字体、颜色、页面、框）去 \texttt{preamble.tex} 改。

\section{LaTeX 基本概念}

\begin{itemize}
  \item \textbf{命令}：反斜杠开头，如 \texttt{\textbackslash section}。
  \item \textbf{参数}：命令可带，放 \texttt{\{\}} 里，如 \texttt{\textbackslash textbf\{粗体\}}。
  \item \textbf{环境}：\texttt{\textbackslash begin\{名\}} 与 \texttt{\textbackslash end\{名\}} 包住一段内容。
  \item \textbf{注释}：\texttt{\%} 及其后整行不编译。
  \item \textbf{段落}：一个空行 = 分段。
\end{itemize}

\section{文字排版}

\subsection{效果}

\textbf{这行加粗。}  \textit{这行斜体（中文变成楷体）。}  \emph{这行强调。}  \texttt{这行等宽字体。}

\subsection{源码}

\begin{codebox}[latex]{文字样式}
\textbf{粗体}          % \textbf
\textit{斜体=楷体}     % \textit（中文自动成楷体）
\emph{强调}            % \emph
\texttt{等宽字体}      % \texttt
\end{codebox}

\begin{notebox}[tl]{中文斜体说明}
中文没有传统意义上的斜体，按国内习惯用\textbf{楷体}代替。这套模板已配好：写 \texttt{\textbackslash textit\{某句\}}，出来就是楷体，无需额外设置。
\end{notebox}

\subsection{标题层级}

\subsubsection*{这一层是 subsubsection}

\subsubsection*{它是怎么写的}

\begin{codebox}[latex]{标题命令}
\chapter{章标题}        % 第一章最顶
\section{节标题}
\subsection{小节标题}
\subsubsection{小小节标题}
\end{codebox}

一个 \texttt{.tex} 文件只能有一个 \texttt{\textbackslash chapter}，下面可以有很多 \texttt{\textbackslash section} 及更小的标题。编号（第零章、0.1、0.1.1）是自动排的。

\section{列表}

\subsection{效果}

无序列表：

\begin{itemize}
  \item 第一项
  \item 第二项
\end{itemize}

有序列表：

\begin{enumerate}
  \item 第一步
  \item 第二步
\end{enumerate}

\subsection{源码}

\begin{codebox}[latex]{列表}
% 无序
\begin{itemize}
  \item 第一项
  \item 第二项
\end{itemize}
% 有序
\begin{enumerate}
  \item 第一步
  \item 第二步
\end{enumerate}
\end{codebox}

\section{数学公式}

\subsection{行内公式}

效果：设 $G(s)$ 是传递函数，$s$ 是复变量，输入是 $u$，输出是 $y$。

源码：\texttt{\textbackslash textbf\{设\} \$G(s)\$ 是传递函数}。

\subsection{独立公式（带编号）}

\begin{equation}
  G(s) \;=\; C\,(sI - A)^{-1}B
  \label{eq:cs}
\end{equation}

\begin{codebox}[latex]{独立公式（带编号）}
\begin{equation}
  G(s) = C\,(sI - A)^{-1}B
  \label{eq:cs}
\end{equation}
\end{codebox}

\texttt{equation} 自动编号，编号是「章.序」（如 0.1、0.2）。

\subsection{独立公式（不带编号）}

\[ G(s) \;=\; \frac{Y(s)}{U(s)} \]

源码：\texttt{\textbackslash[\ G(s) = \textbackslash frac\{Y(s)\}\{U(s)\}\ \textbackslash]}，或用 \texttt{\textbackslash [ ... \textbackslash ]}。

\subsection{多行公式}

\begin{align}
  \dot{x} &= Ax + Bu \\
  y      &= Cx + Du
\end{align}

\begin{codebox}[latex]{多行推导}
\begin{align}
  \dot{x} &= Ax + Bu \\
  y      &= Cx + Du
\end{align}
\end{codebox}

\subsection{常用符号}

这段是渲染出来的：$x^n$　$a_1$　$\frac{a}{b}$　$\sqrt{x}$　$\sum_{i=1}^{n} i$　$\int_0^\infty f(t)\,\mathrm{d}t$　$\alpha\ \beta\ \omega\ \sigma\ \Delta$

\begin{codebox}[latex]{常用符号}
上标 $x^n$     下标 $a_1$
分式 \frac{a}{b}   根式 \sqrt{x}
求和 \sum_{i=1}^{n} i   积分 \int_0^\infty f(t)\,dt
希腊 \alpha \beta \omega \sigma \Delta
\end{codebox}

\subsection{给公式加标签、引用}

效果：式~\eqref{eq:cs} 是状态空间到传递函数的公式。

源码：公式里写 \texttt{\textbackslash label\{名字\}}，别处写 \texttt{\textbackslash eqref\{名字\}}（自带括号）或 \texttt{\textbackslash ref\{名字\}}。

\section{图片}

\subsection{TikZ：直接在模板里画}

这套模板推荐用 TikZ 画图，不依赖外部图片文件、可编译、可随文移动。模板已加载 \texttt{arrows.meta}、\texttt{positioning}、\texttt{calc}、\texttt{shapes.geometric}。

效果（单位反馈闭环框图）：

\begin{center}
\begin{tikzpicture}[>=Stealth, node distance=1.6cm]
  \node[draw, circle, minimum size=0.9cm, inner sep=0pt] (sum) at (0,0) {$+$};
  \node[draw, rectangle, minimum width=2.2cm, minimum height=1.1cm] (plant) at (3.4,0) {$G(s)$};
  \draw[->] (-2.4,0) node[left] {$r$} -- (sum);
  \draw[->] (sum) -- (plant);
  \draw[->] (plant) -- (6.0,0) node[right] {$y$};
  \draw[->] (6.0,0) -- (6.0,-1.6) node[below, midway, fill=white] {$H(s)$} -| (-1.5,-1.6) -- (sum);
\end{tikzpicture}
\end{center}

\begin{codebox}[latex]{上面的框图源码}
\begin{tikzpicture}[>=Stealth, node distance=1.6cm]
  \node[draw, circle, minimum size=0.9cm] (sum) at (0,0) {$+$};
  \node[draw, rectangle, minimum width=2.2cm, minimum height=1.1cm]
        (plant) at (3.4,0) {$G(s)$};
  \draw[->] (-2.4,0) node[left] {$r$} -- (sum);
  \draw[->] (sum) -- (plant);
  \draw[->] (plant) -- (6.0,0) node[right] {$y$};
  \draw[->] (6.0,0) -- (6.0,-1.6)
        node[below, midway, fill=white] {$H(s)$} -| (-1.5,-1.6) -- (sum);
\end{tikzpicture}
\end{codebox}

\begin{notebox}[tl]{节点命名}
每个 \texttt{\textbackslash node} 可起名字（写在 \texttt{()}，如 \texttt{(sum)}、\texttt{(plant)}），后面 \texttt{\textbackslash draw} 直接引用名字连线，比硬写坐标方便。
\end{notebox}

\subsection{TikZ 常用画法（几种够用的）}

\subsubsection{基本图形与文字}

效果：

\begin{center}
\tikz{
  \draw (0,0) rectangle (2,1.4);
  \draw (3,0.7) circle (0.7);
  \draw (4.8,0.7) ellipse (0.8 and 0.55);
  \node at (1,0.7) {矩形};
  \node at (3,0.7) {圆};
  \node at (4.8,0.7) {椭圆};
}
\end{center}

源码：

\begin{codebox}[latex]{基本图形}
\begin{tikzpicture}
  \draw (0,0) rectangle (2,1.4);   % 矩形
  \draw (3,0.7) circle (0.7);        % 圆
  \draw (4.8,0.7) ellipse (0.8 and 0.55); % 椭圆
  \node at (1,0.7) {矩形};          % 在坐标处放文字
  \node at (3,0.7) {圆};
\end{tikzpicture}
\end{codebox}

\subsubsection{相对定位（\texttt{positioning} 库）}

不用手算坐标，用 \texttt{right=of}、\texttt{below=of} 让节点互相挨着。效果：

\begin{center}
\begin{tikzpicture}[node distance=1.2cm]
  \node[draw] (a) {A};
  \node[draw, right=of a] (b) {B};
  \node[draw, below=of b] (c) {C};
  \draw[->] (a) -- (b);
  \draw[->] (b) -- (c);
\end{tikzpicture}
\end{center}

源码：

\begin{codebox}[latex]{相对定位}
\begin{tikzpicture}[node distance=1.2cm]
  \node[draw] (a) {A};
  \node[draw, right=of a] (b) {B};     % 放在 a 右边
  \node[draw, below=of b] (c) {C};     % 放在 b 下方
  \draw[->] (a) -- (b);
  \draw[->] (b) -- (c);
\end{tikzpicture}
\end{codebox}

\subsubsection{箭头样式与颜色}

效果：

\begin{center}
\tikz{
  \draw[->] (0,0) -- (2.2,0);
  \draw[thick,->,blue] (0,0.9) -- (2.2,0.9);
  \draw[dashed,red] (0,1.8) -- (2.2,1.8);
}
\end{center}

源码：

\begin{codebox}[latex]{箭头样式}
\draw[->] (0,0) -- (2.2,0);            % 普通箭头
\draw[thick,->,blue] (0,0.9) -- (2.2,0.9); % 加粗 + 蓝色
\draw[dashed,red] (0,1.8) -- (2.2,1.8);    % 虚线 + 红色
\end{codebox}

\subsubsection{画函数曲线（一阶阶跃响应）}

控制理论常画响应曲线。效果：

\begin{center}
\begin{tikzpicture}[>=Stealth]
  \draw[->] (0,0) -- (5.2,0) node[right]{$t$};
  \draw[->] (0,0) -- (0,3) node[left]{$y$};
  \draw[very thick,blue] plot[smooth,domain=0:4.8] (\x,{2.5*(1-exp(-\x))});
  \draw[dashed] (0,2.5) -- (4.9,2.5) node[above]{$K$};
\end{tikzpicture}
\end{center}

源码：

\begin{codebox}[latex]{函数曲线}
\begin{tikzpicture}[>=Stealth]
  \draw[->] (0,0) -- (5.2,0) node[right]{$t$};      % 横轴
  \draw[->] (0,0) -- (0,3) node[left]{$y$};         % 纵轴
  % plot 画平滑曲线，domain 是 x 范围，\x 是变量
  \draw[very thick,blue] plot[smooth,domain=0:4.8] (\x,{2.5*(1-exp(-\x))});
  \draw[dashed] (0,2.5) -- (4.9,2.5) node[above]{$K$}; % 渐近线
\end{tikzpicture}
\end{codebox}

\subsubsection{坐标轴与标注}

效果：

\begin{center}
\begin{tikzpicture}[>=Stealth]
  \draw[->] (0,0) -- (4,0) node[right]{$\operatorname{Re}$};
  \draw[->] (0,0) -- (0,3) node[left]{$\operatorname{Im}$};
  \fill (1.5,2) circle (2pt) node[above]{$p$};      % 一个点 + 标注
  \draw[dashed] (0,2) -- (1.5,2);
  \draw[dashed] (1.5,0) -- (1.5,2);
\end{tikzpicture}
\end{center}

源码：

\begin{codebox}[latex]{坐标轴与标注}
\begin{tikzpicture}[>=Stealth]
  \draw[->] (0,0) -- (4,0) node[right]{$\operatorname{Re}$};
  \draw[->] (0,0) -- (0,3) node[left]{$\operatorname{Im}$};
  \fill (1.5,2) circle (2pt) node[above]{$p$};       % 点 + 标注
  \draw[dashed] (0,2) -- (1.5,2);                    % 虚线投影
  \draw[dashed] (1.5,0) -- (1.5,2);
\end{tikzpicture}
\end{codebox}

\subsection{外部图片}

先把图片放进 \texttt{images/}，再 \texttt{includegraphics} 插入。

\begin{codebox}[latex]{插入一张图}
\begin{figure}[htbp]
  \centering
  \includegraphics[width=0.6\textwidth]{images/example.png}
  \caption{图的标题}
  \label{fig:example}
\end{figure}
\end{codebox}

\texttt{width=0.6\textbackslash textwidth} 是图片宽度（正文宽度的六成）。\texttt{\textbackslash caption} 写标题，\texttt{\textbackslash label} 给图起名，之后用 \texttt{\textbackslash ref\{名字\}} 引用它。

\section{表格}

\subsection{效果}

\begin{table}[htbp]
  \centering
  \begin{tabular}{lcc}
    \toprule
    名称 & 符号 & 单位 \\
    \midrule
    质量 & $m$ & kg \\
    阻尼 & $c$ & N\,s/m \\
    \bottomrule
  \end{tabular}
  \caption{一个三线表}
\end{table}

\subsection{源码}

\begin{codebox}[latex]{三线表}
\begin{tabular}{lcc}
  \toprule
  名称 & 符号 & 单位 \\
  \midrule
  质量 & $m$ & kg \\
  阻尼 & $c$ & N\,s/m \\
  \bottomrule
\end{tabular}
\end{codebox}

\texttt{lcc} 是列格式（左对齐、居中、居中）。\texttt{\textbackslash toprule}/\texttt{\textbackslash midrule}/\texttt{\textbackslash bottomrule} 是 booktabs 的三条横线。

\section{彩色知识块}

\textbf{所有知识块的样式都在 \texttt{preamble.tex} 配好，你只填内容，不写样式参数。}

\subsection{带编号的知识块}

写法：\texttt{\textbackslash begin\{环境名\}\{标题\}\{标签\}}。第 1 个参数标题（可写一对空 \texttt{\{\}} 只显示编号），第 2 个是供引用的标签。

效果：

\begin{definition}{单边拉普拉斯变换}{def:laplace}
对定义在 $t\ge 0$ 的信号，其拉普拉斯变换为
\begin{equation}
  \mathcal{L}\{f(t)\} = \int_0^\infty f(t)\,e^{-st}\,\mathrm{d}t .
\end{equation}
\end{definition}

\begin{theorem}{}{thm:central}
若系统在原点渐近稳定，则对任意正定 $Q$，方程 $A^{\top}P+PA=-Q$ 有唯一正定解 $P$。
\end{theorem}

\begin{example}{一阶系统阶跃响应}{ex:first}
对 $G(s)=\dfrac{K}{Ts+1}$，单位阶跃响应为 $y(t)=K(1-e^{-t/T})$。
\end{example}

写法：\texttt{\textbackslash begin\{环境名\}\{标题\}\{标签\}} ... \texttt{\textbackslash end\{环境名\}}。上面的定义、定理、例就是这么写的。

现有类型与颜色：

\begin{table}[htbp]
  \centering
  \begin{tabular}{lll}
    \toprule
    环境名 & 显示名 & 颜色 \\
    \midrule
    \texttt{definition} & 定义 & 蓝 \\
    \texttt{theorem}    & 定理 & 红 \\
    \texttt{lemma}      & 引理 & 靛 \\
    \texttt{proposition}& 命题 & 紫 \\
    \texttt{corollary}  & 推论 & 粉 \\
    \texttt{example}    & 例   & 绿 \\
    \texttt{remark}     & 注   & 灰 \\
    \bottomrule
  \end{tabular}
\end{table}

\subsection{不加编号的 \texttt{notebox}}

效果：

\begin{notebox}[og]{证明思路}
先找 $T$ 的不变子空间，再做三角化。
\end{notebox}

\begin{notebox}[pu]{要点}
这是要点框。多行内容之间空一行分段。
\end{notebox}

写法：\texttt{\textbackslash begin\{notebox\}[颜色]\{标题\}} ... \texttt{\textbackslash end\{notebox\}}。方括号选颜色（不写默认蓝），花括号是标题。

可选颜色：\texttt{bl}（蓝）、\texttt{tl}（青）、\texttt{gr}（绿）、\texttt{og}（橙）、\texttt{pu}（紫）、\texttt{rd}（红）、\texttt{in}（靛）、\texttt{pk}（粉）、\texttt{gy}（灰）。

\subsection{自定义一种新颜色、新知识块}

效果、源码：

\begin{codebox}[latex]{新增一种知识块}
% 在 preamble.tex 里：
\definecolor{青}{HTML}{1F9B8E}
\newcoloredtheorem{styling}{风格}{青}{sty}
\end{codebox}

注册后就能 \texttt{\textbackslash begin\{styling\}\{标题\}\{标签\}}。第 3 个参数是颜色名，第 4 个是标签前缀。

\section{插入代码块（\texttt{codebox}）}

\texttt{codebox} 自动加行号、高亮关键字、自动断行。语法：\texttt{\textbackslash begin\{codebox\}\ [语言]\{标题\}}。

效果：

\begin{codebox}[bash]{编译命令}
latexmk -xelatex -synctex=1 -auxdir=build main.tex
\end{codebox}

支持\textbf{多种语言}（\texttt{bash}、\texttt{latex}、\texttt{text}、\texttt{c}、\texttt{python}、\texttt{matlab} 等），每种都会自动高亮\textbf{关键字（蓝）、字符串（绿）}，注释会显示成灰字。

\begin{codebox}[text]{可选语言（方括号里填名字）}
bash              shell 脚本
latex / tex       LaTeX 源码
text              无高亮（目录、说明）
C  c / python / matlab
\end{codebox}

\section{如何新增一个章节}

\begin{enumerate}
  \item 复制 \texttt{chapters/00\_getting\_started.tex} 为 \texttt{chapters/NN\_名字.tex}。
  \item 文件首行保留 \texttt{\% !TEX root = ../main.tex}；把 \texttt{\textbackslash chapter\{名字\}} 改成新名。
  \item 在 \texttt{main.tex} 正文区加一行 \texttt{\textbackslash include\{chapters/NN\_名字\}}。
\end{enumerate}

注意 \texttt{main.tex} 里已经 \texttt{\textbackslash include} 了第零章，你新写的章节在它下面再加一行即可。

% 注：下面 codebox 是「演示写法」，只示意新章节该加在哪一行。真正的
% main.tex 里，第零章那一行（那行 include）已经在了，这里故意不写出来，
% 否则 vimtex 会把教程里这段演示文字误判成「这个文件包含了它自己」而报
% Recursive file inclusion。你写新章节时，照这行格式在 include 列表里加一行即可。

\begin{codebox}[latex]{在 main.tex 里登记（示意新章节加在哪）}
\mainmatter
\include{chapters/01_statespace}   % 你新加的这行
\end{codebox}

\textbf{章节顺序}由 \texttt{\textbackslash include} 的先后决定。

\section{我们配好的宏与环境（重点）}

下面这些都是在 \texttt{preamble.tex} 里配好的\textbf{宏/环境}，你只填内容，\textbf{一个样式参数都不用写}。

\subsection{带编号的彩色框（定义/定理/例…）}

\textbf{效果}：

\begin{proposition}{能控性的秩判据}{prop:rank}
系统 $(A,B)$ 能控，当且仅当能控性矩阵 $\mathcal{C}=[B\; AB\; \dots\; A^{n-1}B]$ 满秩。
\end{proposition}

\begin{corollary}{}{cor:rank}
当 $\mathcal{C}$ 行满秩时，可以通过状态反馈任意配置极点。
\end{corollary}

\begin{codebox}[latex]{这类框的写法}
% 有标题：{标题}{标签}
\begin{proposition}{标题}{标签}  内容 \end{proposition}
% 不写标题：只显示编号
\begin{corollary}{}{标签}  内容 \end{corollary}
\end{codebox}

可用的名字：\texttt{definition}(定义) \texttt{theorem}(定理) \texttt{lemma}(引理) \texttt{proposition}(命题) \texttt{corollary}(推论) \texttt{example}(例) \texttt{remark}(注)。

\subsection{notebox：不加编号的彩框}

\textbf{效果}：

\begin{notebox}[gr]{这一部分解决什么}
适合放「这一段要讲什么」「一个要点」「一段提示」。方括号选颜色，不写默认蓝。
\end{notebox}

\begin{codebox}[latex]{notebox 写法}
\begin{notebox}[gr]{标题}  内容  \end{notebox}
% 不写颜色=蓝：\begin{notebox}{标题}...\end{notebox}
\end{codebox}

\subsection{codebox：代码块}

\textbf{效果}：

\begin{codebox}[latex]{一段 LaTeX 源码}
\begin{equation}\label{eq:demo}
  F(s) = \mathcal{L}\{f(t)\}
\end{equation}
\end{codebox}

写法：\texttt{\textbackslash begin\{codebox\}}\texttt{[语言]}\texttt{\{标题\}} ... \texttt{\textbackslash end\{codebox\}}。第一个参数是语言，第二个是标题。

\subsection{自定义一种新颜色、新框}

\textbf{在 \texttt{preamble.tex} 里加两行}：

\begin{codebox}[latex]{新增一种框}
\definecolor{青}{HTML}{1F9B8E}                 % 1. 定义颜色
\newcoloredtheorem{styling}{风格}{青}{sty}      % 2. 注册成新环境
\end{codebox}

之后就能 \texttt{\textbackslash begin\{styling\}\{标题\}\{标签\}}。\texttt{\textbackslash newcoloredtheorem} 第 1 参数=环境名，第 2=显示名，第 3=颜色名，第 4=标签前缀。


\subsection{颜色代号对照}

\begin{codebox}[text]{可用颜色（notebox 的方括号里填）}
bl=蓝   tl=青   gr=绿   og=橙   pu=紫
rd=红   in=靛   pk=粉   gy=灰
\end{codebox}

\subsection{宏速查表}

\begin{table}[htbp]
  \centering
  \begin{tabular}{lll}
    \toprule
    宏 / 环境 & 作用 & 用法 \\
    \midrule
    \texttt{definition} 等 & 带编号彩框 & \texttt{\textbackslash begin\{环境\}\{标题\}\{标签\}} \\
    \texttt{notebox} & 无编号彩框 & \texttt{\textbackslash begin\{notebox\}[颜色]\{标题\}} \\
    \texttt{codebox} & 代码高亮框 & \texttt{\textbackslash begin\{codebox\}[语言]\{标题\}} \\
    \texttt{\textbackslash newcoloredtheorem} & 新增一种框 & \texttt{（写在 preamble.tex）} \\
    \bottomrule
  \end{tabular}
\end{table}


\section{想改样式、改封面去哪}

\begin{itemize}
  \item 字体 / 配色 / 页面 / 标题字号 / 框样式：都在 \texttt{preamble.tex}。
  \item 封面：\texttt{main.tex} 的 \texttt{\textbackslash begin\{titlepage\}}，TikZ 精确定位，改标题、作者、GitHub 链接都在那。
\end{itemize}

\section{常用命令速查}

这一节把上面用到的命令都列一遍，方便随时翻。分组如下：

\subsection{标题层级}

\texttt{\textbackslash chapter}（章） / \texttt{\textbackslash section}（节） / \texttt{\textbackslash subsection}（小节） / \texttt{\textbackslash subsubsection}（小小节）。

\subsection{文字样式}

\texttt{\textbackslash textbf\{…\}}（加粗） / \texttt{\textbackslash textit\{…\}}（斜体\,=楷体） / \texttt{\textbackslash emph\{…\}}（强调） / \texttt{\textbackslash texttt\{…\}}（等宽）。

\subsection{列表}

\texttt{itemize}（无序） / \texttt{enumerate}（有序），条目用 \texttt{\textbackslash item}。

\subsection{数学}

行内 \texttt{\$…\$}；独立公式 \texttt{equation}（带编号）、\texttt{\textbackslash[}…\texttt{\textbackslash]}（不带编号）；多行 \texttt{align}；\texttt{\textbackslash frac}、\texttt{\textbackslash sqrt}、\texttt{\textbackslash sum}、\texttt{\textbackslash int}；\texttt{\textbackslash alpha} 等希腊字母。

\subsection{插图}

\texttt{includegraphics}（插外部图） / \texttt{tikzpicture}（TikZ 画图）；配 \texttt{figure} 环境、\texttt{\textbackslash caption}、\texttt{\textbackslash label}。

\subsection{表格}

\texttt{tabular}，三线用 \texttt{\textbackslash toprule} / \texttt{\textbackslash midrule} / \texttt{\textbackslash bottomrule}。

\subsection{引用与文件}

\texttt{\textbackslash label}（起名） / \texttt{\textbackslash ref} / \texttt{\textbackslash eqref}（引用）；\texttt{\textbackslash input}（并入） / \texttt{\textbackslash include}（含章节）。

\subsection{本模板的环境}

\texttt{definition}、\texttt{theorem}、\texttt{lemma}、\texttt{proposition}、\texttt{corollary}、\texttt{example}、\texttt{remark}（带编号彩框）；\texttt{notebox}（无编号）；\texttt{codebox}（代码块）；\texttt{\textbackslash newcoloredtheorem}（新增一种）。

\section{用 git / lazygit 管理这份笔记}

这份笔记本身是个 git 仓库，用来记录每次改动、能回退、能备份到 GitHub。你不需要记很多 git 命令，用 \texttt{lazygit} 这个图形化面板就够了。

\subsection{nvim 里分两套工具}

\begin{itemize}
  \item \textbf{lazygit}：全屏操作台，负责\textbf{提交、看历史、切分支、同步远端}。在 nvim 里按 \texttt{\textless leader\textgreater gg} 打开。
  \item \textbf{gitsigns}：行内视图，负责\textbf{看到哪几行改了、单块暂存/回滚、行内 blame}。它贴在编辑器行号旁，不用单独开面板。
\end{itemize}

\subsection{lazygit 的按键}

\begin{table}[htbp]
\centering
\begin{tabular}{lll}
\toprule
按键 & 作用 & 说明 \\
\midrule
\texttt{\textless leader\textgreater gg} & 打开 git 面板 & 提交、看历史、分支都在这 \\
\texttt{\textless leader\textgreater gf} & 当前文件所属仓库 & 定位到本文件所在的 git 项目 \\
\texttt{\textless leader\textgreater gl} & 仓库提交历史 & 看所有提交 \\
\texttt{\textless leader\textgreater gc} & 当前文件提交历史 & 只看这个文件改过啥 \\
\bottomrule
\end{tabular}
\end{table}

\subsection{gitsigns 的按键（改动块）}

\begin{table}[htbp]
\centering
\begin{tabular}{lll}
\toprule
按键 & 作用 & 说明 \\
\midrule
\texttt{]h} / \texttt{[h} & 下一个/上一个改动块 & 跳到你改过的地方 \\
\texttt{\textless leader\textgreater hs} & 暂存当前改动块 & 把这一段加进提交 \\
\texttt{\textless leader\textgreater hr} & 回滚当前改动块 & 撤销这一段 \\
\texttt{\textless leader\textgreater hS} & 暂存整个文件 & \\
\texttt{\textless leader\textgreater hR} & 回滚整个文件 & \\
\texttt{\textless leader\textgreater hp} & 预览改动块 & \\
\texttt{\textless leader\textgreater hb} & 当前行 blame & 看这行谁改的 \\
\texttt{\textless leader\textgreater ht} & 开关行内 blame & \\
\texttt{\textless leader\textgreater hd} & 与索引对比 & \\
\texttt{\textless leader\textgreater hD} & 与上次提交对比 & \\
\bottomrule
\end{tabular}
\end{table}

\subsection{日常提交流程}

\begin{enumerate}
  \item 写完 \texttt{.tex}，按 \texttt{\textbackslash ll} 编译，确认没报错、PDF 正常。
  \item 按 \texttt{\textless leader\textgreater gg} 打开 lazygit，看改动列表。
  \item 在 lazygit 里：回车看差异 → 空格暂存 → \texttt{c} 输入提交信息 → 回车提交。
  \item 需要备份到 GitHub 时，在 lazygit 里按 \texttt{P}（push）推上去。
\end{enumerate}

\begin{notebox}[tl]{记住}
这套模板每改一次就编译一次、看一次效果，攒够一个「有意义」的改动（比如新写了一节）就提交一次。提交信息写清楚改了什么，方便以后查历史。
\end{notebox}

\section{常见的坑}

\begin{notebox}[gy]{编译失败先看这里}
- \texttt{\textbackslash ll} 后按 \texttt{\textbackslash le} 看错误，红色 \texttt{!} 那行即出错位置。
- 中文/特殊字符报错：要用 \texttt{xelatex}（模板默认已用），别用 \texttt{pdflatex}。
- \texttt{\$} 或 \texttt{\{\}} 不配对：这是最常见的编译错误，数一下有没有闭合。
- 图片找不到：确认文件在 \texttt{images/}，路径写对。
- 引用显示 \texttt{??}：没编两遍，多按几次 \texttt{\textbackslash ll}。
\end{notebox}

\section{小结}

文字直接打、要样式用命令；公式用 \texttt{\$} 或 \texttt{equation}；图优先 TikZ、要插图用 \texttt{includegraphics}；知识块只填内容（样式在 preamble 配好）；代码用 \texttt{codebox}；新章节复制文件加 \texttt{\textbackslash include}。出错先看 \texttt{\textbackslash le}。

% 新增章节就在上面加一行，例如：
% \include{chapters/01_statespace}

% ---------- 附录 ----------
\appendix
% \include{chapters/appendix}

\end{document}
